\MathBookSourcePage{346}
\subsection{速度逼近的收敛结果}
\label{sec:ch12-velocity-convergence}
\setcounter{thmcount}{0}
\setcounter{equation}{0}

作为上述理论的第一个应用, 考虑二维 Stokes 方程~\eqref{eq:ch12-stokes-equations} 的逼近空间族. 由~\eqref{eq:ch12-coercivity-on-Z} 定义的条件对形式~\eqref{eq:ch12-stokes-a-form} 在所有 $z\in V$ 上成立, 因而只需证明可逼近性.

令 $V_h^k$ 表示多边形区域 $\Omega\subset\mathbb R^2$ 的一个非退化三角剖分上次数为 $k$ 的 $C^0$ 分片多项式, 其最大三角形直径为 $h$. 令
\MBSourceItem{1}
\begin{equation}
\MathBookFormulaID{formula-ch12-stokes-velocity-space}
\label{eq:ch12-stokes-velocity-space}
V_h:=\{\chvec v\in V_h^k\times V_h^k:\chvec v=0\text{ 于 }\partial\Omega\}
\end{equation}
且令 $\Pi_h$ 是
$\Pi=\{q\in L^2(\Omega):\int_\Omega q(x)\,dx=0\}$ 的任意子集, 满足
\MBSourceItem{2}
\begin{equation}
\MathBookFormulaID{formula-ch12-pressure-approximation-property}
\label{eq:ch12-pressure-approximation-property}
\inf_{q\in\Pi_h}\|p-q\|_{L^2(\Omega)}
\leq Ch^s\|p\|_{H^s(\Omega)}
\qquad\forall p\in\Pi\cap H^s(\Omega),
\end{equation}
其中 $0\leq s\leq k$. 一族这样的压力空间是 Taylor--Hood 压力空间(Brezzi 与 Falk 1991), 如下结果所述.

\MBSourceItem{3}
\begin{Lemma}
\label{lem:ch12-Taylor-Hood-pressure-approximation}
条件~\eqref{eq:ch12-pressure-approximation-property} 对
\MBSourceItem{4}
\begin{equation}
\MathBookFormulaID{formula-ch12-Taylor-Hood-pressure-space}
\label{eq:ch12-Taylor-Hood-pressure-space}
\Pi_h=\left\{q\in V_h^{k-1}:\int_\Omega q(x)\,dx=0\right\}
\end{equation}
成立.
\end{Lemma}

\begin{Proof}
由推论~\MBUnresolvedReference{cor:ch04-source-4-4-24}{4.4.24}(亦参见定理~\MBUnresolvedReference{thm:ch04-source-4-8-12}{4.8.12}),
\[
\MathBookFormulaID{formula-ch12-pressure-projection-error}
\inf_{q\in V_h^{k-1}}\|p-q\|_{L^2(\Omega)}
\leq Ch^s\|p\|_{H^s(\Omega)},\qquad0\leq s\leq k.
\]
下确界由 $p$ 到 $V_h^{k-1}$ 的 $L^2(\Omega)$ 投影 $P_{V_h^{k-1}}p$ 达到. 由于常函数包含在 $V_h^{k-1}$ 中,
$\int_\Omega P_{V_h^{k-1}}p\,dx=\int_\Omega p\,dx=0$, 故 $P_{V_h^{k-1}}p\in\Pi_h$.
\end{Proof}

\MBSourceItem{5}
\begin{Theorem}
\label{thm:ch12-Taylor-Hood-velocity-error}
设 $V_h$ 如~\eqref{eq:ch12-stokes-velocity-space} 给出, $\Pi_h$ 满足~\eqref{eq:ch12-pressure-approximation-property}. 令 $u,u_h$ 如定理~\ref{thm:ch12-discrete-velocity-error}. 假设 $k\geq4$. 则对 $0\leq s\leq k$,
\[
\MathBookFormulaID{formula-ch12-Taylor-Hood-velocity-error}
\|u-u_h\|_{H^1(\Omega)^2}
\leq Ch^s\left(\|u\|_{H^{s+1}(\Omega)^2}
+\|p\|_{H^s(\Omega)}\right),
\]
前提是~\eqref{eq:ch12-stokes-abstract-problem} 的解 $(u,p)$ 满足
$(u,p)\in H^{s+1}(\Omega)^2\times H^s(\Omega)$.
\end{Theorem}

\begin{Proof}
由~\eqref{eq:ch12-pressure-approximation-property} 和定理~\ref{thm:ch12-discrete-velocity-error}, 只需证明
\MBSourceItem{6}
\begin{equation}
\MathBookFormulaID{formula-ch12-Zh-velocity-approximation}
\label{eq:ch12-Zh-velocity-approximation}
\inf_{v\in Z_h}\|u-v\|_{H^1(\Omega)^2}
\leq Ch^s\|u\|_{H^{s+1}(\Omega)^2}.
\end{equation}
注意 $u\in Z$ 蕴含 $u=\chvec{\operatorname{curl}}\psi$, 其中
$\psi\in H^{s+2}(\Omega)\cap\mathring H^2(\Omega)$(Arnold, Scott 与 Vogelius 1988). 使用 $\mathcal T^h$ 上的 Argyris 元, 存在分片多项式 $\psi_h\in\mathring H^2(\Omega)$, 次数为 $s+1$, 使得
\[
\MathBookFormulaID{formula-ch12-Argyris-approximation}
\|\psi-\psi_h\|_{H^2(\Omega)}
\leq Ch^s\|\psi\|_{H^{s+2}(\Omega)}.
\]
由于 $\chvec{\operatorname{curl}}\psi_h\in Z_h$, 证明完成. 又因 $u=\chvec{\operatorname{curl}}\psi$, 有
$\|\psi\|_{H^{s+2}(\Omega)}\leq C\|u\|_{H^{s+1}(\Omega)^2}$.
\end{Proof}

\MBSourceItem{7}
\begin{Remark}
\label{rem:ch12-k3-velocity}
下一节证明 $k=2$ 的结果. $k=3$ 的情形由 Brezzi 与 Falk (1991) 处理.
\end{Remark}

\MathBookSourcePage{347}
为用混合方法逼近标量椭圆问题~\eqref{eq:ch12-porous-medium-pressure-equation}, 必须考虑相应形式 $a(\cdot,\cdot)$ 并不在整个 $V$ 上强制. 它显然在空间
\[
\MathBookFormulaID{formula-ch12-scalar-elliptic-kernel}
Z=\{\chvec v\in H(\operatorname{div}):\operatorname{div}\chvec v=0\}
\]
上强制, 故~\eqref{eq:ch12-reduced-kernel-problem} 适定. 然而, 必须保证它在~\eqref{eq:ch12-discrete-kernel-Zh} 给出的 $Z_h$ 上也适定. 一个简单办法是保证 $Z_h\subset Z$, 下面将说明一种实现方式.

回到上一节的一般记号, 注意 $\Pi_h$ 自然地与 $DV_h$ 配对. 若取 $\Pi_h=DV_h$, 则~\eqref{eq:ch12-discrete-kernel-Zh} 对 $Z_h$ 的定义保证 $Z_h\subset Z$, 从而保证强制性. 例如, 可取 $V_h=V_h^k\times V_h^k$(参见~\eqref{eq:ch12-stokes-velocity-space}), 于是~\eqref{eq:ch12-Zh-velocity-approximation} 的证明表明
\MBSourceItem{8}
\begin{equation}
\MathBookFormulaID{formula-ch12-scalar-elliptic-kernel-approximation}
\label{eq:ch12-scalar-elliptic-kernel-approximation}
\inf_{v\in Z_h}\|u-v\|_{H^1(\Omega)^2}
\leq Ch^k\|u\|_{H^{k+1}(\Omega)^2},
\end{equation}
因为该结果在以 $V_h\cap\{v\in H^1(\Omega)^2:\operatorname{div}v=0\}$ 代替 $Z_h$ 时成立, 而后者是 $Z_h$ 的子集.

作为~\eqref{eq:ch12-pressure-approximation-property} 和~\eqref{eq:ch12-scalar-elliptic-kernel-approximation} 的推论, 得到以下结果.

\MBSourceItem{9}
\begin{Theorem}
\label{thm:ch12-scalar-elliptic-mixed-velocity}
考虑标量椭圆问题~\eqref{eq:ch12-porous-medium-pressure-equation} 的混合方法, $n=2$, 形式 $a(\cdot,\cdot)$ 由~\eqref{eq:ch12-porous-medium-a-form} 给出. 令 $V_h=V_h^k\times V_h^k$, $\Pi_h=\operatorname{div}V_h$. 假设 $k\geq4$. 则对 $0\leq s\leq k$,
\[
\MathBookFormulaID{formula-ch12-scalar-elliptic-mixed-velocity}
\|u-u_h\|_{H^1(\Omega)^2}
\leq Ch^{s+1}\|\chvec{\operatorname{grad}}p\|_{H^{s+1}(\Omega)^2}
+Ch^s\|p\|_{H^s(\Omega)},
\]
前提是~\eqref{eq:ch12-porous-medium-pressure-equation} 的解 $p$ 满足
$\chvec{\operatorname{grad}}p\in H^{s+1}(\Omega)^2$.
\end{Theorem}

类似地, 由~\eqref{eq:ch12-Cea-discrete-kernel} 与~\eqref{eq:ch12-Zh-velocity-approximation} 得到:

\MBSourceItem{10}
\begin{Theorem}
\label{thm:ch12-Taylor-Hood-velocity-divVh}
令 $V_h$ 如~\eqref{eq:ch12-stokes-velocity-space} 给出, 并令 $\Pi_h=\operatorname{div}V_h$. 令 $u,u_h$ 如定理~\ref{thm:ch12-discrete-velocity-error}. 假设 $k\geq4$. 则对 $0\leq s\leq k$,
\[
\MathBookFormulaID{formula-ch12-Taylor-Hood-velocity-divVh}
\|u-u_h\|_{H^1(\Omega)^2}
\leq Ch^s\|u\|_{H^{s+1}(\Omega)^2},
\]
前提是~\eqref{eq:ch12-stokes-abstract-problem} 的解 $u$ 满足
$u\in H^{s+1}(\Omega)^2$.
\end{Theorem}

\MathBookSourcePage{348}
上述理论也可用于 Stokes 方程的非协调有限元逼近. 定义
\[
\MathBookFormulaID{formula-ch12-nonconforming-stokes-forms}
a_h(u,v):=\sum_{T\in\mathcal T^h}\int_T
\chten{\operatorname{grad}}_h u:\chten{\operatorname{grad}}_h v\,dx,
\qquad
b_h(v,q):=-\sum_{T\in\mathcal T^h}\int_T\operatorname{div}_h v\,q\,dx,
\]
并令 $\|v\|_h:=\sqrt{a_h(v,v)}$. 令 $V_h^*$ 为~\eqref{eq:ch11-nonconforming-vector-space} 定义的空间, 并令 $\Pi_h$ 表示 $\mathcal T^h$ 上的分片常函数 $q$, 满足 $\int_\Omega q\,dx=0$. 由
\MBSourceItem{11}
\begin{equation}
\MathBookFormulaID{formula-ch12-nonconforming-stokes-problem}
\label{eq:ch12-nonconforming-stokes-problem}
\begin{aligned}
a_h(u_h,v)+b_h(v,p_h)&=\int_\Omega f\cdot v\,dx&&\forall v\in V_h^*,\\
b_h(u_h,q)&=0&&\forall q\in\Pi_h
\end{aligned}
\end{equation}
定义 $u_h\in V_h^*$ 和 $p_h\in\Pi_h$.

\MBSourceItem{12}
\begin{Theorem}
\label{thm:ch12-nonconforming-stokes-velocity}
若 $u\in H^2(\Omega)^2$ 且 $p\in H^1(\Omega)$, 则
\[
\MathBookFormulaID{formula-ch12-nonconforming-stokes-velocity}
\|u-u_h\|_h\leq Ch\left(\|u\|_{H^2(\Omega)^2}
+\|p\|_{H^1(\Omega)}\right).
\]
\end{Theorem}

\begin{Proof}
令 $Z_h=\{v\in V_h^*:b_h(v,q)=0\ \forall q\in\Pi_h\}$. 由非协调 Strang 估计,
\[
\MathBookFormulaID{formula-ch12-nonconforming-stokes-strang}
\|u-u_h\|_h\leq\inf_{v\in Z_h}\|u-v\|_h+
\sup_{w\in Z_h\setminus\{0\}}
\frac{|a_h(u-u_h,w)|}{\|w\|_h}.
\]
$a_h(\cdot,\cdot)$ 在 $V=\mathring H^1(\Omega)^2$ 上强制, 在非退化 $V_h^*$ 上也是如此, 因为 $a_h(v,v)=0$ 蕴含 $v$ 是分片常数, 且零边界条件与中点连续性蕴含 $v=0$. 然后
\[
\MathBookFormulaID{formula-ch12-nonconforming-stokes-residual}
\begin{aligned}
a_h(u-u_h,w)
&=a_h(u,w)-\int_\Omega f\cdot w\,dx\\
&=a_h(u,w)-\int_\Omega(-\Delta u+\chvec{\operatorname{grad}}p)\cdot w\,dx.
\end{aligned}
\]
分部积分并约定 $w$ 在 $\Omega$ 外为零, 从引理~\MBUnresolvedReference{lem:ch10-source-3-7}{10.3.7} 和~\MBUnresolvedReference{lem:ch10-source-3-9}{10.3.9} 得到
\[
\MathBookFormulaID{formula-ch12-nonconforming-stokes-consistency-bound}
|a_h(u-u_h,w)|
\leq Ch\left(\|u\|_{H^2(\Omega)^2}+\|p\|_{H^1(\Omega)}\right)\|w\|_h.
\]
沿定理~\ref{thm:ch12-discrete-velocity-error} 的推导,
$|b_h(w,p)|\leq Ch\|w\|_h\|p\|_{H^1(\Omega)}$. 只需证明
$\inf_{v\in Z_h}\|u-v\|_h\leq Ch\|u\|_{H^2(\Omega)^2}$. 令 $M_h$ 为图~\MBUnresolvedReference{fig:ch10-source-10-5}{10.5} 所示 Morley 空间. 由 $\chvec{\operatorname{curl}}M_h\subset Z_h$ 及定理~\ref{thm:ch12-Taylor-Hood-velocity-error} 的证明完成.
\end{Proof}
